% f03_claim1 variant c -- the square that closes.
% Data: the structural argument for Proposition (mapping/topology
%       orthogonality) in main.tex, sec:abi-orth -- the three edge annotations
%       are its three sentences, naming build_schedule, generate_instructions
%       and the F-sb-gate hypothesis of Claim 1. The root word is the N=128 row
%       of tables/invariance.tex (0x415F7E8FF9581062); the empirical closure
%       counts are \WallTBs{} (data/coverage_macros.tex) and \InvNonSettling{}
%       (data/results_macros.tex).
% Strategy: argue FORMALLY rather than by exhibit. a and b show three fabrics
% and invite the reader to compare roots; this one shows that the comparison
% never had to be made -- the right-hand edge of the square is an identity by
% construction, and each of the other three edges carries the reason it is.
% Colour: every box on the compliant construction path is spBlue; the two
% observed roots and the closing identity are the confirmed invariant (spGreen);
% the edge reasons are annotation, so they stay spMute and take no series
% colour. Box weight and position carry the same hierarchy in greyscale.
% Target: IEEE double column, 7.16in = 18.2cm. Drawn inside 16.4cm.
\begin{tikzpicture}[
  x=1cm, y=1cm,
  bx/.style={draw=spBlue, fill=spBlueF, line width=0.45pt, align=center,
             font=\tiny, inner sep=3pt, minimum height=0.70cm, rounded corners=1.5pt,
             text=spInk},
  rb/.style={draw=spGreen, fill=spGreenF, line width=0.8pt, align=center,
             font=\tiny, inner sep=3pt, minimum height=0.70cm, rounded corners=1.5pt,
             text=spInk},
  ar/.style={draw=spBlue, line width=0.5pt, -{Latex[length=1.7mm,width=1.2mm]}},
  lbl/.style={font=\tiny, align=center, inner sep=1.5pt, text=spInk},
  note/.style={font=\tiny, align=left, inner sep=1.5pt, text=spMute},
]
% ---------------- the fixed data, above the square -------------------------
\node[bx, text width=3.5cm] (SRC) at (2.10,3.95)
  {fixed: $N$, leaf vector $L$,\\assignment, refinement, FADD unit};
\node[bx, text width=3.5cm] (SCH) at (2.10,2.55)
  {per-PE schedule $s$\\ \sv{build\_schedule}};
\draw[ar] (SRC) -- (SCH);
\node[note, anchor=west] at (3.95,3.25)
  {a \texttt{COMPUTE}'s $(\texttt{addr\_a},\texttt{addr\_b},\texttt{dest})$\\
   depends only on the assignment and the\\
   DMEM allocator --- neither reads $\mu$ or $\rho$};

% ---------------- the square ----------------------------------------------
\node[bx, text width=3.7cm] (TL) at (2.10,1.05)
  {image on $(\mu_1,\rho_1)$: ring\\ \sv{generate\_instructions}};
\node[bx, text width=3.7cm] (BL) at (2.10,-0.55)
  {image on $(\mu_2,\rho_2)$: fat-tree\\ \sv{generate\_instructions}};
\node[rb, text width=3.5cm] (TR) at (9.90,1.05)
  {$R_\mathrm{observed}(\mu_1,\rho_1)$\\ \texttt{0x415F7E8FF9581062}};
\node[rb, text width=3.5cm] (BR) at (9.90,-0.55)
  {$R_\mathrm{observed}(\mu_2,\rho_2)$\\ \texttt{0x415F7E8FF9581062}};
\draw[ar] (SCH) -- (TL);
\draw[ar, draw=spGreen] (TL) -- node[lbl, above, text=spGreen] {execute} (TR);
\draw[ar, draw=spGreen] (BL) -- node[lbl, below, text=spGreen] {execute} (BR);
\draw[ar] (TL) -- node[lbl, left, text=spBlue, align=right]
  {re-map,\\re-route} (BL);
% the closing edge
\draw[spGreen, line width=0.8pt] (11.70,1.05) -- (12.15,1.05)
  -- (12.15,-0.55) -- (11.70,-0.55);
\node[font=\small, text=spGreen] at (12.60,0.25) {$=$};
\node[note, anchor=west, text=spGreen] at (12.90,0.25)
  {$\Rcan(N,L)$};

% ---------------- edge reasons --------------------------------------------
\node[note, anchor=north west] at (4.05,0.72)
  {a \texttt{SEND}'s $(\texttt{dest\_gpu},\texttt{target\_addr},\text{link})$
   \emph{does} depend on $\mu$ and $\rho$ ---\\
   but it fixes only \emph{where} data flow, never the operand binding
   at the receiver};
\node[note, anchor=north west] at (4.05,-0.92)
  {\textsf{F-sb-gate}: the receiving \texttt{COMPUTE} reads DMEM by
   \emph{address},\\ not by arrival time, port or virtual channel};

% ---------------- the reading ---------------------------------------------
\node[lbl, anchor=north, text width=15.6cm, align=center, text=spInk] at (7.20,-1.95)
  {The two axes that change most obviously under cluster reshape --- the
   mapping $\mu$ and the routing tables $\rho$ --- are structurally orthogonal
   to the bit result, so the square closes by inspection of the code path that
   emits the instructions rather than by an empirical average. The campaign is
   consistent with it: \WallTBs{} settled simulations, \InvNonSettling{}
   non-settling, and no finite root differing from $\Rcan(N,L)$.};
\end{tikzpicture}
